\documentclass[10pt,a4paper]{amsart}
\usepackage[margin=19mm]{geometry}
\usepackage{amsmath,amssymb,mathtools}
\usepackage[scr=rsfs,cal=euler]{mathalfa}
\usepackage{xfrac}
\usepackage[unicode,hidelinks,bookmarks=false]{hyperref}
\newcommand{\ie}{i.e.\ }
\newcommand{\cf}{cf.\ }
\newcommand{\resp}{respectively\ }
\newcommand{\Subsection}{\subsection}
\providecommand{\nobreakdash}{\nobreak\hbox{-}}
\setlength{\emergencystretch}{2em}
\multlinegap0pt
\usepackage{amssymb}
\usepackage{xcolor}
\newtheorem{proposition}{Proposition}
\newtheorem{lemma}{Lemma}
\newtheorem{theorem}{Theorem}
\title{Finsler geometry in anisotropic superconductivity: a Ginzburg-Landau approach}
\author{Y. Alipour Fakhri}
\date{Source: arXiv:2510.19720v4 (CC BY 4.0)}
\begin{document}
\maketitle

\noindent\emph{Source and licence.} Finsler geometry in anisotropic superconductivity: a Ginzburg-Landau approach. Version arXiv:2510.19720v4 (CC BY 4.0), distributed by arXiv under the Creative Commons Attribution 4.0 International licence, http://creativecommons.org/licenses/by/4.0/, as recorded in the arXiv licence metadata for this exact version and shown on the abstract page https://arxiv.org/abs/2510.19720v4. Full licence notice: \texttt{licenses/arxiv-2510.19720-CC-BY-4.0.txt}.\par\smallskip

\noindent\textit{Editorial scope.} Counted unit: the article's theorem thm:EL-Dirichlet (source0.tex lines 391--400). The article's complete original proof of it is delivered with it (source0.tex lines 402--421, one \texttt{proof} environment of 20 lines). No mathematics is written, repaired or re-derived: the statement and the proof are the article's own lines, in the article's own notation, and no retained line is substituted or re-flowed.  Retained uncounted as the source-local context the proof needs: lines 296--310; lines 360--366.  Nothing was omitted from any retained span, so the package carries no omission record.  No retained line contains a citation, so the wrapper carries no bibliography and no bibliography entry.  No retained line contains a figure environment, an image-inclusion command or drawing code, so no image is delivered and the package declares no asset dependency.  Editorial formatting only, in addition: an AMS \texttt{amsart} preamble that restates the article's own declarations and macros used by the retained lines, namely the environments \texttt{proposition}, \texttt{lemma}, \texttt{theorem} on separate counters, together with the standard packages those lines need.  The retained environments print in this document as Proposition 1, Lemma 1, Theorem 1 in order of appearance, whereas the article prints the same items with the numbering of its own full document.  The complete native source of the article as distributed in the arXiv e-print for this exact version is bundled unchanged as \texttt{sources/187594/source0.tex}.
\begin{proposition}\label{prop:Legendre-diffeo}
%[Legendre diffeomorphism and basic identities]
For each $x\in M$, $L_x:T_xM\setminus\{0\}\to T_x^*M\setminus\{0\}$ is a $C^\infty$ diffeomorphism.
Its inverse is given by $L_x^{-1}(\xi)=\partial_\xi\!\big(\tfrac 12 F^*(x,\xi)^2\big)$.
Moreover, for all $y\neq 0$ and $\xi\neq 0$,
\begin{equation}\label{eq:F-Fstar-compat}
F^*\!\big(x,L_x(y)\big)\ =\ F(x,y),\qquad
F\!\big(x,L_x^{-1}(\xi)\big)\ =\ F^*(x,\xi),
\end{equation}
and the Fenchel--Young relation holds:
\begin{equation}\label{eq:Fenchel-Young}
\frac 12 F(x,y)^2 + \frac 12 F^*(x,\xi)^2 \ \ge\ \xi(y).
\end{equation}
Equality hold if and only if $\xi=L_x(y)$, equivalently, $y=L_x^{-1}(\xi)$.
\end{proposition}

\begin{lemma}\label{lem:ibp}
%[Integration by parts]
For $u\in C^\infty(M)$ and $X\in C^1(TM)$,
\begin{equation}\label{eq:ibp}
\int_M \! du(X)\,d\mu_F\ =\ -\int_M \! u\ \mathrm{div}_{\mu_F}X\,d\mu_F.
\end{equation}
\end{lemma}

\begin{theorem}\label{thm:EL-Dirichlet}
%[Euler--Lagrange equation for $\mathcal{E}_F$]
For $u\in C^\infty(M)$ and $\varphi\in C^\infty(M)$,
\begin{align*}
\frac{d}{dt}\,\mathcal{E}_F[u+t\varphi]\Big|_{t=0}
\ =\ -\int_M \varphi\, \Delta^F_{\mu_F}u\, d\mu_F.
\end{align*}
Equivalently, the critical points of $\mathcal{E}_F$ are the (weak) solutions of
$\Delta^F_{\mu_F} u=0$.
\end{theorem}

\begin{proof}
Set $\xi_t:=d(u+t\varphi)=du+t\,d\varphi$. Using Proposition \ref{prop:Legendre-diffeo},
\begin{align*}
\frac{d}{dt}\,\frac 12 F^*(x,\xi_t)^2\Big|_{t=0}
\ =\ \Big\langle \partial_\xi\!\Big(\tfrac 12 F^*(x,\xi)^2\Big)\Big|_{\xi=du}\!,\, d\varphi\Big\rangle
\ =\ \langle \nabla_F u,\ d\varphi\rangle,
\end{align*}
where the last pairing is the natural one $T_xM\times T_x^*M\to\mathbb{R}$. Hence
\begin{align*}
\frac{d}{dt}\,\mathcal{E}_F[u+t\varphi]\Big|_{t=0}
\ =\ \int_M \! \langle \nabla_F u,\ d\varphi\rangle\, d\mu_F
\ =\ \int_M \! d\varphi(\nabla_F u)\, d\mu_F.
\end{align*}
Now apply Lemma \ref{lem:ibp} with $X=\nabla_F u$ to obtain
\begin{align*}
\int_M d\varphi(\nabla_F u)\, d\mu_F
\ =\ -\int_M \varphi\, \mathrm{div}_{\mu_F}(\nabla_F u)\, d\mu_F
\ =\ -\int_M \varphi\, \Delta^F_{\mu_F}u\, d\mu_F.
\end{align*}
\end{proof}

\end{document}
